
\documentclass{article}
\usepackage{colortbl}
\usepackage{makecell}
\usepackage{multirow}
\usepackage{supertabular}

\begin{document}

\newcounter{utterance}

\centering \large Interaction Transcript for game `clean\_up', experiment `0\_easy\_5obj\_bg', episode 0 with deepseek{-}v4{-}pro.
\vspace{24pt}

{ \footnotesize  \setcounter{utterance}{1}
\setlength{\tabcolsep}{0pt}
\begin{supertabular}{c@{$\;$}|p{.15\linewidth}@{}p{.15\linewidth}p{.15\linewidth}p{.15\linewidth}p{.15\linewidth}p{.15\linewidth}}
    \# & $\;$A & \multicolumn{4}{c}{Game Master} & $\;\:$B\\
    \hline

    \theutterance \stepcounter{utterance}  
    & & \multicolumn{4}{p{0.6\linewidth}}{
        \cellcolor[rgb]{0.9,0.9,0.9}{
            \makecell[{{p{\linewidth}}}]{
                \texttt{\tiny{[P1$\langle$GM]}}
                \texttt{Аз съм вашият ръководител на играта и вие играете съвместна игра със следната мрежа като игрално поле:} \\
\\ 
\texttt{```} \\
\texttt{ 1234567} \\
\texttt{╔═╤═══╤═╗} \\
\texttt{║◌│◌◌◌│◌║ 1} \\
\texttt{╟─┼──┬┘◌║ 2} \\
\texttt{║◌│◌◌│◌◌║ 3} \\
\texttt{╟─┴──┘H◌║ 4} \\
\texttt{║◌◌◌◌◌I◌║ 5} \\
\texttt{║C◌◌◌◌◌◌║ 6} \\
\texttt{║◌W◌◌T◌◌║ 7} \\
\texttt{╚═══════╝} \\
\\ 
\texttt{```} \\
\\ 
\texttt{* Горният ръб показва x{-}координати, които нарастват надясно, а десният ръб показва y{-}координати, които нарастват надолу.} \\
\texttt{* Следните обекти са разположени на случаен принцип върху вашата мрежа: 'W', 'I', 'T', 'C', 'H'.} \\
\\ 
\\ 
\texttt{Другият играч вижда вариант на игралното поле, където обектите са разположени на различни произволни места. Вие не можете да видите полето на другия играч, а той не може да види вашето.} \\
\\ 
\\ 
\texttt{**Цел:**} \\
\\ 
\texttt{И двамата играчи трябва да преместят обектите на своя фон така, че еднакви обекти да се окажат на едни и същи координати. Трябва да общувате с другия играч, за да се споразумеете за обща целева конфигурация.} \\
\\ 
\\ 
\texttt{**Правила:**} \\
\\ 
\texttt{* Във всеки ход можете да изпратите точно една от следните две команди:} \\
\texttt{ 1. `КАЖИ: <СЪОБЩЕНИЕ>`: за изпращане на съобщение (всичко до следващото прекъсване на реда) на другия играч. Аз ще го препратя на вашия партньор.} \\
\texttt{2. `ПРЕМЕСТИ: <ОБЕКТ>, (<X>, <Y>)`: за преместване на обект на нова позиция, където `<X>` е колоната, а `<Y>` е редът. Ще ви информирам дали преместването ви е успешно или не.} \\
\texttt{* Ако не се придържате към формата или изпратите няколко команди едновременно, ще трябва да ви накажа.} \\
\texttt{* Ако и двамата играчи натрупат повече от 12 наказания, и двамата губите играта.} \\
\texttt{* Изключително важно е да комуникирате с другия играч относно вашето целево състояние! Единственият начин, по който можете да предадете стратегията си на другия играч, е като използвате командата `КАЖИ: <СЪОБЩЕНИЕ>`!} \\
\\ 
\\ 
\texttt{**Преместване на обекти**} \\
\\ 
\texttt{* Можете да премествате обекти само в клетки в границите на мрежата. Целевата клетка трябва да е празна, т.е. да съдържа само символа '◌'.} \\
\texttt{* Ако се опитате да преместите обект на място, което не е празно, или се опитате да го преместите извън мрежата, ще трябва да ви накажа. Получавате още един опит.} \\
\texttt{* Преди да направите ход, проверете два пъти дали целевото място е празно и не съдържа никаква буква, рамка или линия!} \\
\\ 
\\ 
\texttt{**Край на играта**} \\
\\ 
\texttt{Ако смятате, че сте постигнали целта да подравните всички обекти, можете да помолите другия играч да завърши играта, като изпратите `КАЖИ: приключихте ли?`. Ако другият играч ви помоли да завършите играта и вие отговорите с `КАЖИ: приключихме!`, играта ще приключи.} \\
\\ 
\texttt{И двамата играчи печелят, ако играта приключи в рамките на 20 рунда, където един рунд се определя като двама играчи, всеки от които изпраща валидна команда.} \\
\\ 
\texttt{**Оценяване:**} \\
\\ 
\texttt{Колкото по{-}близо са еднаквите обекти в двете игрални полета, толкова повече точки получавате. Наказанията намаляват точките ви. Можете ли да подобрите рекорда?} \\
\\ 
\\ 
\texttt{Моля, изпратете съобщение на другия играч, за да започнете играта!} \\
            }
        }
    }
    & \\ \\

    \theutterance \stepcounter{utterance}  
    & \multicolumn{4}{p{0.6\linewidth}}{
        \cellcolor[rgb]{0.9,0.9,0.9}{
            \makecell[{{p{\linewidth}}}]{
                \texttt{\tiny{[P1$\rangle$GM]}}
                \texttt{КАЖИ: Здравей! Да координираме позициите. Моите обекти са на: W(2,7), I(6,5), T(5,7), C(1,6), H(5,4). Къде са твоите?} \\
            }
        }
    }
    & & \\ \\

    \theutterance \stepcounter{utterance}  
    & & & \multicolumn{2}{p{0.3\linewidth}}{
        \cellcolor[rgb]{0.9,0.9,0.9}{
            \makecell[{{p{\linewidth}}}]{
                \texttt{\tiny{[GM$|$GM]}}
                \texttt{КАЖИ: Здравей! Да координираме позициите. Моите обекти са на: W(2,7), I(6,5), T(5,7), C(1,6), H(5,4). Къде са твоите?} \\
            }
        }
    }
    & & \\ \\

    \theutterance \stepcounter{utterance}  
    & & \multicolumn{4}{p{0.6\linewidth}}{
        \cellcolor[rgb]{0.9,0.9,0.9}{
            \makecell[{{p{\linewidth}}}]{
                \texttt{\tiny{[GM$\rangle$P2]}}
                \texttt{Аз съм вашият ръководител на играта и вие играете съвместна игра със следната мрежа като игрално поле:} \\
\\ 
\texttt{```} \\
\texttt{ 1234567} \\
\texttt{╔═╤═══╤═╗} \\
\texttt{║◌│◌◌◌│H║ 1} \\
\texttt{╟─┼──┬┘◌║ 2} \\
\texttt{║◌│◌◌│◌C║ 3} \\
\texttt{╟─┴──┘◌T║ 4} \\
\texttt{║◌◌◌◌◌◌◌║ 5} \\
\texttt{║◌◌◌I◌◌◌║ 6} \\
\texttt{║◌◌◌◌W◌◌║ 7} \\
\texttt{╚═══════╝} \\
\\ 
\texttt{```} \\
\\ 
\texttt{* Горният ръб показва x{-}координати, които нарастват надясно, а десният ръб показва y{-}координати, които нарастват надолу.} \\
\texttt{* Следните обекти са разположени на случаен принцип върху вашата мрежа: 'W', 'I', 'T', 'C', 'H'.} \\
\\ 
\\ 
\texttt{Другият играч вижда вариант на игралното поле, където обектите са разположени на различни произволни места. Вие не можете да видите полето на другия играч, а той не може да види вашето.} \\
\\ 
\\ 
\texttt{**Цел:**} \\
\\ 
\texttt{И двамата играчи трябва да преместят обектите на своя фон така, че еднакви обекти да се окажат на едни и същи координати. Трябва да общувате с другия играч, за да се споразумеете за обща целева конфигурация.} \\
\\ 
\\ 
\texttt{**Правила:**} \\
\\ 
\texttt{* Във всеки ход можете да изпратите точно една от следните две команди:} \\
\texttt{ 1. `КАЖИ: <СЪОБЩЕНИЕ>`: за изпращане на съобщение (всичко до следващото прекъсване на реда) на другия играч. Аз ще го препратя на вашия партньор.} \\
\texttt{2. `ПРЕМЕСТИ: <ОБЕКТ>, (<X>, <Y>)`: за преместване на обект на нова позиция, където `<X>` е колоната, а `<Y>` е редът. Ще ви информирам дали преместването ви е успешно или не.} \\
\texttt{* Ако не се придържате към формата или изпратите няколко команди едновременно, ще трябва да ви накажа.} \\
\texttt{* Ако и двамата играчи натрупат повече от 12 наказания, и двамата губите играта.} \\
\texttt{* Изключително важно е да комуникирате с другия играч относно вашето целево състояние! Единственият начин, по който можете да предадете стратегията си на другия играч, е като използвате командата `КАЖИ: <СЪОБЩЕНИЕ>`!} \\
\\ 
\\ 
\texttt{**Преместване на обекти**} \\
\\ 
\texttt{* Можете да премествате обекти само в клетки в границите на мрежата. Целевата клетка трябва да е празна, т.е. да съдържа само символа '◌'.} \\
\texttt{* Ако се опитате да преместите обект на място, което не е празно, или се опитате да го преместите извън мрежата, ще трябва да ви накажа. Получавате още един опит.} \\
\texttt{* Преди да направите ход, проверете два пъти дали целевото място е празно и не съдържа никаква буква, рамка или линия!} \\
\\ 
\\ 
\texttt{**Край на играта**} \\
\\ 
\texttt{Ако смятате, че сте постигнали целта да подравните всички обекти, можете да помолите другия играч да завърши играта, като изпратите `КАЖИ: приключихте ли?`. Ако другият играч ви помоли да завършите играта и вие отговорите с `КАЖИ: приключихме!`, играта ще приключи.} \\
\\ 
\texttt{И двамата играчи печелят, ако играта приключи в рамките на 20 рунда, където един рунд се определя като двама играчи, всеки от които изпраща валидна команда.} \\
\\ 
\texttt{**Оценяване:**} \\
\\ 
\texttt{Колкото по{-}близо са еднаквите обекти в двете игрални полета, толкова повече точки получавате. Наказанията намаляват точките ви. Можете ли да подобрите рекорда?} \\
\\ 
\texttt{Другият играч започна играта, като изпрати това съобщение:} \\
\\ 
\texttt{"Здравей! Да координираме позициите. Моите обекти са на: W(2,7), I(6,5), T(5,7), C(1,6), H(5,4). Къде са твоите?"} \\
\\ 
\texttt{Каква е първата ви команда?} \\
            }
        }
    }
    & \\ \\

    \theutterance \stepcounter{utterance}  
    & & & \multicolumn{4}{p{0.6\linewidth}}{
        \cellcolor[rgb]{0.9,0.9,0.9}{
            \makecell[{{p{\linewidth}}}]{
                \texttt{\tiny{[GM$\langle$P2]}}
                \texttt{КАЖИ: Здравей! Моите обекти са на: W(5,7), I(4,6), T(6,4), C(7,3), H(7,1). Да се разберем за общи позиции {-} може ли всеки обект да отиде там, където е средната позиция между нашите две?} \\
            }
        }
    }
     \\ \\

    \theutterance \stepcounter{utterance}  
    & & & \multicolumn{2}{p{0.3\linewidth}}{
        \cellcolor[rgb]{0.9,0.9,0.9}{
            \makecell[{{p{\linewidth}}}]{
                \texttt{\tiny{[GM$|$GM]}}
                \texttt{КАЖИ: Здравей! Моите обекти са на: W(5,7), I(4,6), T(6,4), C(7,3), H(7,1). Да се разберем за общи позиции {-} може ли всеки обект да отиде там, където е средната позиция между нашите две?} \\
            }
        }
    }
    & & \\ \\

    \theutterance \stepcounter{utterance}  
    & & \multicolumn{4}{p{0.6\linewidth}}{
        \cellcolor[rgb]{0.9,0.9,0.9}{
            \makecell[{{p{\linewidth}}}]{
                \texttt{\tiny{[P1$\langle$GM]}}
                \texttt{Вашето съобщение беше предадено на другия играч.} \\
\\ 
\\ 
\texttt{В момента играете рунд 1 от максимум 20.} \\
\texttt{Колективно сте натрупали 0 от 12 наказания.} \\
\texttt{Другият играч изпрати това съобщение:} \\
\\ 
\\ 
\texttt{"Здравей! Моите обекти са на: W(5,7), I(4,6), T(6,4), C(7,3), H(7,1). Да се разберем за общи позиции {-} може ли всеки обект да отиде там, където е средната позиция между нашите две?"} \\
\\ 
\\ 
\texttt{Каква е вашата следваща команда?} \\
            }
        }
    }
    & \\ \\

    \theutterance \stepcounter{utterance}  
    & \multicolumn{4}{p{0.6\linewidth}}{
        \cellcolor[rgb]{0.9,0.9,0.9}{
            \makecell[{{p{\linewidth}}}]{
                \texttt{\tiny{[P1$\rangle$GM]}}
                \texttt{Нека изчисля средните позиции между нашите обекти:} \\
\\ 
\texttt{**W:** Аз (2,7) | Ти (5,7) → средно (3.5, 7) → нека закръглим до **(3,7)** или **(4,7)**. Предлагам **(3,7)**.} \\
\\ 
\texttt{**I:** Аз (6,5) | Ти (4,6) → средно (5, 5.5) → **(5,5)** или **(5,6)**. Предлагам **(5,5)**.} \\
\\ 
\texttt{**T:** Аз (5,7) | Ти (6,4) → средно (5.5, 5.5) → **(5,5)** или **(6,5)** или **(5,6)**. Но (5,5) вече е заето от I! Предлагам **(6,5)** за T.} \\
\\ 
\texttt{**C:** Аз (1,6) | Ти (7,3) → средно (4, 4.5) → **(4,4)** или **(4,5)**. Предлагам **(4,4)**.} \\
\\ 
\texttt{**H:** Аз (5,4) | Ти (7,1) → средно (6, 2.5) → **(6,2)** или **(6,3)**. Предлагам **(6,2)**.} \\
\\ 
\texttt{КАЖИ: Добра идея! Предлагам тези целеви позиции: W(3,7), I(5,5), T(6,5), C(4,4), H(6,2). Провери дали тези клетки са празни в твоето поле. Съгласен ли си?} \\
            }
        }
    }
    & & \\ \\

    \theutterance \stepcounter{utterance}  
    & & & \multicolumn{2}{p{0.3\linewidth}}{
        \cellcolor[rgb]{0.9,0.9,0.9}{
            \makecell[{{p{\linewidth}}}]{
                \texttt{\tiny{[GM$|$GM]}}
                \texttt{Нека изчисля средните позиции между нашите обекти:} \\
\\ 
\texttt{**W:** Аз (2,7) | Ти (5,7) → средно (3.5, 7) → нека закръглим до **(3,7)** или **(4,7)**. Предлагам **(3,7)**.} \\
\\ 
\texttt{**I:** Аз (6,5) | Ти (4,6) → средно (5, 5.5) → **(5,5)** или **(5,6)**. Предлагам **(5,5)**.} \\
\\ 
\texttt{**T:** Аз (5,7) | Ти (6,4) → средно (5.5, 5.5) → **(5,5)** или **(6,5)** или **(5,6)**. Но (5,5) вече е заето от I! Предлагам **(6,5)** за T.} \\
\\ 
\texttt{**C:** Аз (1,6) | Ти (7,3) → средно (4, 4.5) → **(4,4)** или **(4,5)**. Предлагам **(4,4)**.} \\
\\ 
\texttt{**H:** Аз (5,4) | Ти (7,1) → средно (6, 2.5) → **(6,2)** или **(6,3)**. Предлагам **(6,2)**.} \\
\\ 
\texttt{КАЖИ: Добра идея! Предлагам тези целеви позиции: W(3,7), I(5,5), T(6,5), C(4,4), H(6,2). Провери дали тези клетки са празни в твоето поле. Съгласен ли си?} \\
            }
        }
    }
    & & \\ \\

    \theutterance \stepcounter{utterance}  
    & & & \multicolumn{2}{p{0.3\linewidth}}{
        \cellcolor[rgb]{0.9,0.9,0.9}{
            \makecell[{{p{\linewidth}}}]{
                \texttt{\tiny{[GM$|$GM]}}
                \texttt{Invalid format: head is not empty: 'Нека изчисля средните позиции между нашите обекти:} \\
\\ 
\texttt{**W:** Аз (2,7) | Ти (5,7) → средно (3.5, 7) → нека закръглим до **(3,7)** или **(4,7)**. Предлагам **(3,7)**.} \\
\\ 
\texttt{**I:** Аз (6,5) | Ти (4,6) → средно (5, 5.5) → **(5,5)** или **(5,6)**. Предлагам **(5,5)**.} \\
\\ 
\texttt{**T:** Аз (5,7) | Ти (6,4) → средно (5.5, 5.5) → **(5,5)** или **(6,5)** или **(5,6)**. Но (5,5) вече е заето от I! Предлагам **(6,5)** за T.} \\
\\ 
\texttt{**C:** Аз (1,6) | Ти (7,3) → средно (4, 4.5) → **(4,4)** или **(4,5)**. Предлагам **(4,4)**.} \\
\\ 
\texttt{**H:** Аз (5,4) | Ти (7,1) → средно (6, 2.5) → **(6,2)** или **(6,3)**. Предлагам **(6,2)**.} \\
\\ 
\texttt{'} \\
            }
        }
    }
    & & \\ \\

    \theutterance \stepcounter{utterance}  
    & & \multicolumn{4}{p{0.6\linewidth}}{
        \cellcolor[rgb]{0.9,0.9,0.9}{
            \makecell[{{p{\linewidth}}}]{
                \texttt{\tiny{[P1$\langle$GM]}}
                \texttt{Наказание: Вашето съобщение не трябва да съдържа нищо преди командата!} \\
\\ 
\texttt{Уверете се, че вашият отговор съдържа само `КАЖИ: <СЪОБЩЕНИЕ>` или `ПРЕМЕСТИ: <ОБЕКТ>, (<X>, <Y>)` и нищо друго!} \\
\\ 
\texttt{Колективно сте натрупали 1 от 12 наказания.} \\
\texttt{Моля, опитайте отново!} \\
            }
        }
    }
    & \\ \\

    \theutterance \stepcounter{utterance}  
    & \multicolumn{4}{p{0.6\linewidth}}{
        \cellcolor[rgb]{0.9,0.9,0.9}{
            \makecell[{{p{\linewidth}}}]{
                \texttt{\tiny{[P1$\rangle$GM]}}
                \texttt{КАЖИ: Добра идея! Предлагам тези целеви позиции: W(3,7), I(5,5), T(6,5), C(4,4), H(6,2). Провери дали тези клетки са празни в твоето поле. Съгласен ли си?} \\
            }
        }
    }
    & & \\ \\

    \theutterance \stepcounter{utterance}  
    & & & \multicolumn{2}{p{0.3\linewidth}}{
        \cellcolor[rgb]{0.9,0.9,0.9}{
            \makecell[{{p{\linewidth}}}]{
                \texttt{\tiny{[GM$|$GM]}}
                \texttt{КАЖИ: Добра идея! Предлагам тези целеви позиции: W(3,7), I(5,5), T(6,5), C(4,4), H(6,2). Провери дали тези клетки са празни в твоето поле. Съгласен ли си?} \\
            }
        }
    }
    & & \\ \\

    \theutterance \stepcounter{utterance}  
    & & \multicolumn{4}{p{0.6\linewidth}}{
        \cellcolor[rgb]{0.9,0.9,0.9}{
            \makecell[{{p{\linewidth}}}]{
                \texttt{\tiny{[GM$\rangle$P2]}}
                \texttt{Вашето съобщение беше предадено на другия играч.} \\
\\ 
\\ 
\texttt{В момента играете рунд 2 от максимум 20.} \\
\texttt{Колективно сте натрупали 1 от 12 наказания.} \\
\texttt{Другият играч изпрати това съобщение:} \\
\\ 
\\ 
\texttt{"Добра идея! Предлагам тези целеви позиции: W(3,7), I(5,5), T(6,5), C(4,4), H(6,2). Провери дали тези клетки са празни в твоето поле. Съгласен ли си?"} \\
\\ 
\\ 
\texttt{Каква е вашата следваща команда?} \\
            }
        }
    }
    & \\ \\

    \theutterance \stepcounter{utterance}  
    & & & \multicolumn{4}{p{0.6\linewidth}}{
        \cellcolor[rgb]{0.9,0.9,0.9}{
            \makecell[{{p{\linewidth}}}]{
                \texttt{\tiny{[GM$\langle$P2]}}
                \texttt{Нека първо проверя дали предложените целеви позиции са празни в моето поле:} \\
\\ 
\texttt{Моето поле:} \\
\texttt{{-} W(5,7) → целева (3,7): (3,7) е '◌'? Да, празна е.} \\
\texttt{{-} I(4,6) → целева (5,5): (5,5) е '◌'? Да, празна е.} \\
\texttt{{-} T(6,4) → целева (6,5): (6,5) е '◌'? Да, празна е.} \\
\texttt{{-} C(7,3) → целева (4,4): (4,4) е '◌'? Да, празна е.} \\
\texttt{{-} H(7,1) → целева (6,2): (6,2) е '◌'? Да, празна е.} \\
\\ 
\texttt{Всички целеви клетки са празни в моето поле. Съгласен съм.} \\
\\ 
\texttt{КАЖИ: Всички целеви позиции са празни в моето поле. Съгласен съм с W(3,7), I(5,5), T(6,5), C(4,4), H(6,2). Да започваме преместването!} \\
            }
        }
    }
     \\ \\

    \theutterance \stepcounter{utterance}  
    & & & \multicolumn{2}{p{0.3\linewidth}}{
        \cellcolor[rgb]{0.9,0.9,0.9}{
            \makecell[{{p{\linewidth}}}]{
                \texttt{\tiny{[GM$|$GM]}}
                \texttt{Нека първо проверя дали предложените целеви позиции са празни в моето поле:} \\
\\ 
\texttt{Моето поле:} \\
\texttt{{-} W(5,7) → целева (3,7): (3,7) е '◌'? Да, празна е.} \\
\texttt{{-} I(4,6) → целева (5,5): (5,5) е '◌'? Да, празна е.} \\
\texttt{{-} T(6,4) → целева (6,5): (6,5) е '◌'? Да, празна е.} \\
\texttt{{-} C(7,3) → целева (4,4): (4,4) е '◌'? Да, празна е.} \\
\texttt{{-} H(7,1) → целева (6,2): (6,2) е '◌'? Да, празна е.} \\
\\ 
\texttt{Всички целеви клетки са празни в моето поле. Съгласен съм.} \\
\\ 
\texttt{КАЖИ: Всички целеви позиции са празни в моето поле. Съгласен съм с W(3,7), I(5,5), T(6,5), C(4,4), H(6,2). Да започваме преместването!} \\
            }
        }
    }
    & & \\ \\

    \theutterance \stepcounter{utterance}  
    & & & \multicolumn{2}{p{0.3\linewidth}}{
        \cellcolor[rgb]{0.9,0.9,0.9}{
            \makecell[{{p{\linewidth}}}]{
                \texttt{\tiny{[GM$|$GM]}}
                \texttt{Invalid format: head is not empty: 'Нека първо проверя дали предложените целеви позиции са празни в моето поле:} \\
\\ 
\texttt{Моето поле:} \\
\texttt{{-} W(5,7) → целева (3,7): (3,7) е '◌'? Да, празна е.} \\
\texttt{{-} I(4,6) → целева (5,5): (5,5) е '◌'? Да, празна е.} \\
\texttt{{-} T(6,4) → целева (6,5): (6,5) е '◌'? Да, празна е.} \\
\texttt{{-} C(7,3) → целева (4,4): (4,4) е '◌'? Да, празна е.} \\
\texttt{{-} H(7,1) → целева (6,2): (6,2) е '◌'? Да, празна е.} \\
\\ 
\texttt{Всички целеви клетки са празни в моето поле. Съгласен съм.} \\
\\ 
\texttt{'} \\
            }
        }
    }
    & & \\ \\

    \theutterance \stepcounter{utterance}  
    & & \multicolumn{4}{p{0.6\linewidth}}{
        \cellcolor[rgb]{0.9,0.9,0.9}{
            \makecell[{{p{\linewidth}}}]{
                \texttt{\tiny{[GM$\rangle$P2]}}
                \texttt{Наказание: Вашето съобщение не трябва да съдържа нищо преди командата!} \\
\\ 
\texttt{Уверете се, че вашият отговор съдържа само `КАЖИ: <СЪОБЩЕНИЕ>` или `ПРЕМЕСТИ: <ОБЕКТ>, (<X>, <Y>)` и нищо друго!} \\
\\ 
\texttt{Колективно сте натрупали 2 от 12 наказания.} \\
\texttt{Моля, опитайте отново!} \\
            }
        }
    }
    & \\ \\

    \theutterance \stepcounter{utterance}  
    & & & \multicolumn{4}{p{0.6\linewidth}}{
        \cellcolor[rgb]{0.9,0.9,0.9}{
            \makecell[{{p{\linewidth}}}]{
                \texttt{\tiny{[GM$\langle$P2]}}
                \texttt{КАЖИ: Всички целеви позиции са празни в моето поле. Съгласен съм с W(3,7), I(5,5), T(6,5), C(4,4), H(6,2). Да започваме преместването!} \\
            }
        }
    }
     \\ \\

    \theutterance \stepcounter{utterance}  
    & & & \multicolumn{2}{p{0.3\linewidth}}{
        \cellcolor[rgb]{0.9,0.9,0.9}{
            \makecell[{{p{\linewidth}}}]{
                \texttt{\tiny{[GM$|$GM]}}
                \texttt{КАЖИ: Всички целеви позиции са празни в моето поле. Съгласен съм с W(3,7), I(5,5), T(6,5), C(4,4), H(6,2). Да започваме преместването!} \\
            }
        }
    }
    & & \\ \\

    \theutterance \stepcounter{utterance}  
    & & \multicolumn{4}{p{0.6\linewidth}}{
        \cellcolor[rgb]{0.9,0.9,0.9}{
            \makecell[{{p{\linewidth}}}]{
                \texttt{\tiny{[P1$\langle$GM]}}
                \texttt{Вашето съобщение беше предадено на другия играч.} \\
\\ 
\\ 
\texttt{В момента играете рунд 2 от максимум 20.} \\
\texttt{Колективно сте натрупали 2 от 12 наказания.} \\
\texttt{Другият играч изпрати това съобщение:} \\
\\ 
\\ 
\texttt{"Всички целеви позиции са празни в моето поле. Съгласен съм с W(3,7), I(5,5), T(6,5), C(4,4), H(6,2). Да започваме преместването!"} \\
\\ 
\\ 
\texttt{Каква е вашата следваща команда?} \\
            }
        }
    }
    & \\ \\

    \theutterance \stepcounter{utterance}  
    & \multicolumn{4}{p{0.6\linewidth}}{
        \cellcolor[rgb]{0.9,0.9,0.9}{
            \makecell[{{p{\linewidth}}}]{
                \texttt{\tiny{[P1$\rangle$GM]}}
                \texttt{Първо ще проверя кои целеви позиции са празни в моето поле и кои обекти трябва да преместя:} \\
\\ 
\texttt{Моето поле:} \\
\texttt{{-} W(2,7) → цел (3,7): трябва да проверя дали (3,7) е празно. На моето поле (3,7) е ◌? Да, изглежда празно.} \\
\texttt{{-} I(6,5) → цел (5,5): (5,5) е празно? На ред 5: ◌◌◌◌◌I◌, значи (5,5) е празно (◌).} \\
\texttt{{-} T(5,7) → цел (6,5): (6,5) е заето от I в момента! Трябва първо да преместя I.} \\
\texttt{{-} C(1,6) → цел (4,4): (4,4) е H? Не, на ред 4: ╟─┴──┘H◌, значи (4,4) е H. Трябва да преместя H първо.} \\
\texttt{{-} H(5,4) → цел (6,2): (6,2) е ◌? На ред 2: ╟─┼──┬┘◌║, (6,2) е празно.} \\
\\ 
\texttt{Нека започна с местене на I от (6,5) на (5,5).} \\
\\ 
\texttt{ПРЕМЕСТИ: I, (5,5)} \\
            }
        }
    }
    & & \\ \\

    \theutterance \stepcounter{utterance}  
    & & & \multicolumn{2}{p{0.3\linewidth}}{
        \cellcolor[rgb]{0.9,0.9,0.9}{
            \makecell[{{p{\linewidth}}}]{
                \texttt{\tiny{[GM$|$GM]}}
                \texttt{Първо ще проверя кои целеви позиции са празни в моето поле и кои обекти трябва да преместя:} \\
\\ 
\texttt{Моето поле:} \\
\texttt{{-} W(2,7) → цел (3,7): трябва да проверя дали (3,7) е празно. На моето поле (3,7) е ◌? Да, изглежда празно.} \\
\texttt{{-} I(6,5) → цел (5,5): (5,5) е празно? На ред 5: ◌◌◌◌◌I◌, значи (5,5) е празно (◌).} \\
\texttt{{-} T(5,7) → цел (6,5): (6,5) е заето от I в момента! Трябва първо да преместя I.} \\
\texttt{{-} C(1,6) → цел (4,4): (4,4) е H? Не, на ред 4: ╟─┴──┘H◌, значи (4,4) е H. Трябва да преместя H първо.} \\
\texttt{{-} H(5,4) → цел (6,2): (6,2) е ◌? На ред 2: ╟─┼──┬┘◌║, (6,2) е празно.} \\
\\ 
\texttt{Нека започна с местене на I от (6,5) на (5,5).} \\
\\ 
\texttt{ПРЕМЕСТИ: I, (5,5)} \\
            }
        }
    }
    & & \\ \\

    \theutterance \stepcounter{utterance}  
    & & & \multicolumn{2}{p{0.3\linewidth}}{
        \cellcolor[rgb]{0.9,0.9,0.9}{
            \makecell[{{p{\linewidth}}}]{
                \texttt{\tiny{[GM$|$GM]}}
                \texttt{Invalid format: head is not empty: 'Първо ще проверя кои целеви позиции са празни в моето поле и кои обекти трябва да преместя:} \\
\\ 
\texttt{Моето поле:} \\
\texttt{{-} W(2,7) → цел (3,7): трябва да проверя дали (3,7) е празно. На моето поле (3,7) е ◌? Да, изглежда празно.} \\
\texttt{{-} I(6,5) → цел (5,5): (5,5) е празно? На ред 5: ◌◌◌◌◌I◌, значи (5,5) е празно (◌).} \\
\texttt{{-} T(5,7) → цел (6,5): (6,5) е заето от I в момента! Трябва първо да преместя I.} \\
\texttt{{-} C(1,6) → цел (4,4): (4,4) е H? Не, на ред 4: ╟─┴──┘H◌, значи (4,4) е H. Трябва да преместя H първо.} \\
\texttt{{-} H(5,4) → цел (6,2): (6,2) е ◌? На ред 2: ╟─┼──┬┘◌║, (6,2) е празно.} \\
\\ 
\texttt{Нека започна с местене на I от (6,5) на (5,5).} \\
\\ 
\texttt{'} \\
            }
        }
    }
    & & \\ \\

    \theutterance \stepcounter{utterance}  
    & & \multicolumn{4}{p{0.6\linewidth}}{
        \cellcolor[rgb]{0.9,0.9,0.9}{
            \makecell[{{p{\linewidth}}}]{
                \texttt{\tiny{[P1$\langle$GM]}}
                \texttt{Наказание: Вашето съобщение не трябва да съдържа нищо преди командата!} \\
\\ 
\texttt{Уверете се, че вашият отговор съдържа само `КАЖИ: <СЪОБЩЕНИЕ>` или `ПРЕМЕСТИ: <ОБЕКТ>, (<X>, <Y>)` и нищо друго!} \\
\\ 
\texttt{Колективно сте натрупали 3 от 12 наказания.} \\
\texttt{Моля, опитайте отново!} \\
            }
        }
    }
    & \\ \\

    \theutterance \stepcounter{utterance}  
    & \multicolumn{4}{p{0.6\linewidth}}{
        \cellcolor[rgb]{0.9,0.9,0.9}{
            \makecell[{{p{\linewidth}}}]{
                \texttt{\tiny{[P1$\rangle$GM]}}
                \texttt{ПРЕМЕСТИ: I, (5,5)} \\
            }
        }
    }
    & & \\ \\

    \theutterance \stepcounter{utterance}  
    & & & \multicolumn{2}{p{0.3\linewidth}}{
        \cellcolor[rgb]{0.9,0.9,0.9}{
            \makecell[{{p{\linewidth}}}]{
                \texttt{\tiny{[GM$|$GM]}}
                \texttt{ПРЕМЕСТИ: I, (5,5)} \\
            }
        }
    }
    & & \\ \\

    \theutterance \stepcounter{utterance}  
    & & & \multicolumn{2}{p{0.3\linewidth}}{
        \cellcolor[rgb]{0.9,0.9,0.9}{
            \makecell[{{p{\linewidth}}}]{
                \texttt{\tiny{[GM$|$GM]}}
                \texttt{I moved to (5, 5)} \\
            }
        }
    }
    & & \\ \\

    \theutterance \stepcounter{utterance}  
    & & \multicolumn{4}{p{0.6\linewidth}}{
        \cellcolor[rgb]{0.9,0.9,0.9}{
            \makecell[{{p{\linewidth}}}]{
                \texttt{\tiny{[GM$\rangle$P2]}}
                \texttt{Вашето съобщение беше предадено на другия играч.} \\
\\ 
\\ 
\texttt{В момента играете рунд 3 от максимум 20.} \\
\texttt{Колективно сте натрупали 3 от 12 наказания.} \\
\texttt{Другият играч премести обект на своята мрежа.} \\
\\ 
\\ 
\texttt{Каква е вашата следваща команда?} \\
            }
        }
    }
    & \\ \\

    \theutterance \stepcounter{utterance}  
    & & & \multicolumn{4}{p{0.6\linewidth}}{
        \cellcolor[rgb]{0.9,0.9,0.9}{
            \makecell[{{p{\linewidth}}}]{
                \texttt{\tiny{[GM$\langle$P2]}}
                \texttt{Ще започна да местя обектите си към договорените целеви позиции. Нека първо преместя H от (7,1) към (6,2).} \\
\\ 
\texttt{ПРЕМЕСТИ: H, (6,2)} \\
            }
        }
    }
     \\ \\

    \theutterance \stepcounter{utterance}  
    & & & \multicolumn{2}{p{0.3\linewidth}}{
        \cellcolor[rgb]{0.9,0.9,0.9}{
            \makecell[{{p{\linewidth}}}]{
                \texttt{\tiny{[GM$|$GM]}}
                \texttt{Ще започна да местя обектите си към договорените целеви позиции. Нека първо преместя H от (7,1) към (6,2).} \\
\\ 
\texttt{ПРЕМЕСТИ: H, (6,2)} \\
            }
        }
    }
    & & \\ \\

    \theutterance \stepcounter{utterance}  
    & & & \multicolumn{2}{p{0.3\linewidth}}{
        \cellcolor[rgb]{0.9,0.9,0.9}{
            \makecell[{{p{\linewidth}}}]{
                \texttt{\tiny{[GM$|$GM]}}
                \texttt{Invalid format: head is not empty: 'Ще започна да местя обектите си към договорените целеви позиции. Нека първо преместя H от (7,1) към (6,2).} \\
\\ 
\texttt{'} \\
            }
        }
    }
    & & \\ \\

    \theutterance \stepcounter{utterance}  
    & & \multicolumn{4}{p{0.6\linewidth}}{
        \cellcolor[rgb]{0.9,0.9,0.9}{
            \makecell[{{p{\linewidth}}}]{
                \texttt{\tiny{[GM$\rangle$P2]}}
                \texttt{Наказание: Вашето съобщение не трябва да съдържа нищо преди командата!} \\
\\ 
\texttt{Уверете се, че вашият отговор съдържа само `КАЖИ: <СЪОБЩЕНИЕ>` или `ПРЕМЕСТИ: <ОБЕКТ>, (<X>, <Y>)` и нищо друго!} \\
\\ 
\texttt{Колективно сте натрупали 4 от 12 наказания.} \\
\texttt{Моля, опитайте отново!} \\
            }
        }
    }
    & \\ \\

    \theutterance \stepcounter{utterance}  
    & & & \multicolumn{4}{p{0.6\linewidth}}{
        \cellcolor[rgb]{0.9,0.9,0.9}{
            \makecell[{{p{\linewidth}}}]{
                \texttt{\tiny{[GM$\langle$P2]}}
                \texttt{ПРЕМЕСТИ: H, (6,2)} \\
            }
        }
    }
     \\ \\

    \theutterance \stepcounter{utterance}  
    & & & \multicolumn{2}{p{0.3\linewidth}}{
        \cellcolor[rgb]{0.9,0.9,0.9}{
            \makecell[{{p{\linewidth}}}]{
                \texttt{\tiny{[GM$|$GM]}}
                \texttt{ПРЕМЕСТИ: H, (6,2)} \\
            }
        }
    }
    & & \\ \\

    \theutterance \stepcounter{utterance}  
    & & & \multicolumn{2}{p{0.3\linewidth}}{
        \cellcolor[rgb]{0.9,0.9,0.9}{
            \makecell[{{p{\linewidth}}}]{
                \texttt{\tiny{[GM$|$GM]}}
                \texttt{Наказание: (6,2) не е празно, но съдържа '┘'.} \\
\\ 
\texttt{Колективно сте натрупали 5 от 12 наказания.} \\
\texttt{Моля, опитайте отново!} \\
            }
        }
    }
    & & \\ \\

    \theutterance \stepcounter{utterance}  
    & & \multicolumn{4}{p{0.6\linewidth}}{
        \cellcolor[rgb]{0.9,0.9,0.9}{
            \makecell[{{p{\linewidth}}}]{
                \texttt{\tiny{[GM$\rangle$P2]}}
                \texttt{Наказание: (6,2) не е празно, но съдържа '┘'.} \\
\\ 
\texttt{Колективно сте натрупали 5 от 12 наказания.} \\
\texttt{Моля, опитайте отново!} \\
            }
        }
    }
    & \\ \\

    \theutterance \stepcounter{utterance}  
    & & & \multicolumn{4}{p{0.6\linewidth}}{
        \cellcolor[rgb]{0.9,0.9,0.9}{
            \makecell[{{p{\linewidth}}}]{
                \texttt{\tiny{[GM$\langle$P2]}}
                \texttt{ПРЕМЕСТИ: H, (6,1)} \\
            }
        }
    }
     \\ \\

    \theutterance \stepcounter{utterance}  
    & & & \multicolumn{2}{p{0.3\linewidth}}{
        \cellcolor[rgb]{0.9,0.9,0.9}{
            \makecell[{{p{\linewidth}}}]{
                \texttt{\tiny{[GM$|$GM]}}
                \texttt{ПРЕМЕСТИ: H, (6,1)} \\
            }
        }
    }
    & & \\ \\

    \theutterance \stepcounter{utterance}  
    & & & \multicolumn{2}{p{0.3\linewidth}}{
        \cellcolor[rgb]{0.9,0.9,0.9}{
            \makecell[{{p{\linewidth}}}]{
                \texttt{\tiny{[GM$|$GM]}}
                \texttt{Наказание: (6,1) не е празно, но съдържа '│'.} \\
\\ 
\texttt{Колективно сте натрупали 6 от 12 наказания.} \\
\texttt{Моля, опитайте отново!} \\
            }
        }
    }
    & & \\ \\

    \theutterance \stepcounter{utterance}  
    & & \multicolumn{4}{p{0.6\linewidth}}{
        \cellcolor[rgb]{0.9,0.9,0.9}{
            \makecell[{{p{\linewidth}}}]{
                \texttt{\tiny{[GM$\rangle$P2]}}
                \texttt{Наказание: (6,1) не е празно, но съдържа '│'.} \\
\\ 
\texttt{Колективно сте натрупали 6 от 12 наказания.} \\
\texttt{Моля, опитайте отново!} \\
            }
        }
    }
    & \\ \\

    \theutterance \stepcounter{utterance}  
    & & & \multicolumn{4}{p{0.6\linewidth}}{
        \cellcolor[rgb]{0.9,0.9,0.9}{
            \makecell[{{p{\linewidth}}}]{
                \texttt{\tiny{[GM$\langle$P2]}}
                \texttt{ПРЕМЕСТИ: H, (7,2)} \\
            }
        }
    }
     \\ \\

    \theutterance \stepcounter{utterance}  
    & & & \multicolumn{2}{p{0.3\linewidth}}{
        \cellcolor[rgb]{0.9,0.9,0.9}{
            \makecell[{{p{\linewidth}}}]{
                \texttt{\tiny{[GM$|$GM]}}
                \texttt{ПРЕМЕСТИ: H, (7,2)} \\
            }
        }
    }
    & & \\ \\

    \theutterance \stepcounter{utterance}  
    & & & \multicolumn{2}{p{0.3\linewidth}}{
        \cellcolor[rgb]{0.9,0.9,0.9}{
            \makecell[{{p{\linewidth}}}]{
                \texttt{\tiny{[GM$|$GM]}}
                \texttt{H moved to (7, 2)} \\
            }
        }
    }
    & & \\ \\

    \theutterance \stepcounter{utterance}  
    & & \multicolumn{4}{p{0.6\linewidth}}{
        \cellcolor[rgb]{0.9,0.9,0.9}{
            \makecell[{{p{\linewidth}}}]{
                \texttt{\tiny{[P1$\langle$GM]}}
                \texttt{Преместихте 'I' на (5,5) успешно. Вашата актуализирана решетка изглежда така:} \\
\texttt{```} \\
\texttt{ 1234567} \\
\texttt{╔═╤═══╤═╗} \\
\texttt{║◌│◌◌◌│◌║ 1} \\
\texttt{╟─┼──┬┘◌║ 2} \\
\texttt{║◌│◌◌│◌◌║ 3} \\
\texttt{╟─┴──┘H◌║ 4} \\
\texttt{║◌◌◌◌I◌◌║ 5} \\
\texttt{║C◌◌◌◌◌◌║ 6} \\
\texttt{║◌W◌◌T◌◌║ 7} \\
\texttt{╚═══════╝} \\
\texttt{```} \\
\texttt{В момента играете рунд 3 от максимум 20.} \\
\texttt{Колективно сте натрупали 6 от 12 наказания.} \\
\texttt{Другият играч премести обект на своята мрежа.} \\
\\ 
\\ 
\texttt{Каква е вашата следваща команда?} \\
            }
        }
    }
    & \\ \\

    \theutterance \stepcounter{utterance}  
    & \multicolumn{4}{p{0.6\linewidth}}{
        \cellcolor[rgb]{0.9,0.9,0.9}{
            \makecell[{{p{\linewidth}}}]{
                \texttt{\tiny{[P1$\rangle$GM]}}
                \texttt{ПРЕМЕСТИ: T, (6,5)} \\
            }
        }
    }
    & & \\ \\

    \theutterance \stepcounter{utterance}  
    & & & \multicolumn{2}{p{0.3\linewidth}}{
        \cellcolor[rgb]{0.9,0.9,0.9}{
            \makecell[{{p{\linewidth}}}]{
                \texttt{\tiny{[GM$|$GM]}}
                \texttt{ПРЕМЕСТИ: T, (6,5)} \\
            }
        }
    }
    & & \\ \\

    \theutterance \stepcounter{utterance}  
    & & & \multicolumn{2}{p{0.3\linewidth}}{
        \cellcolor[rgb]{0.9,0.9,0.9}{
            \makecell[{{p{\linewidth}}}]{
                \texttt{\tiny{[GM$|$GM]}}
                \texttt{T moved to (6, 5)} \\
            }
        }
    }
    & & \\ \\

    \theutterance \stepcounter{utterance}  
    & & \multicolumn{4}{p{0.6\linewidth}}{
        \cellcolor[rgb]{0.9,0.9,0.9}{
            \makecell[{{p{\linewidth}}}]{
                \texttt{\tiny{[GM$\rangle$P2]}}
                \texttt{Преместихте 'H' на (7,2) успешно. Вашата актуализирана решетка изглежда така:} \\
\texttt{```} \\
\texttt{ 1234567} \\
\texttt{╔═╤═══╤═╗} \\
\texttt{║◌│◌◌◌│◌║ 1} \\
\texttt{╟─┼──┬┘H║ 2} \\
\texttt{║◌│◌◌│◌C║ 3} \\
\texttt{╟─┴──┘◌T║ 4} \\
\texttt{║◌◌◌◌◌◌◌║ 5} \\
\texttt{║◌◌◌I◌◌◌║ 6} \\
\texttt{║◌◌◌◌W◌◌║ 7} \\
\texttt{╚═══════╝} \\
\texttt{```} \\
\texttt{В момента играете рунд 4 от максимум 20.} \\
\texttt{Колективно сте натрупали 6 от 12 наказания.} \\
\texttt{Другият играч премести обект на своята мрежа.} \\
\\ 
\\ 
\texttt{Каква е вашата следваща команда?} \\
            }
        }
    }
    & \\ \\

    \theutterance \stepcounter{utterance}  
    & & & \multicolumn{4}{p{0.6\linewidth}}{
        \cellcolor[rgb]{0.9,0.9,0.9}{
            \makecell[{{p{\linewidth}}}]{
                \texttt{\tiny{[GM$\langle$P2]}}
                \texttt{ПРЕМЕСТИ: C, (4,4)} \\
            }
        }
    }
     \\ \\

    \theutterance \stepcounter{utterance}  
    & & & \multicolumn{2}{p{0.3\linewidth}}{
        \cellcolor[rgb]{0.9,0.9,0.9}{
            \makecell[{{p{\linewidth}}}]{
                \texttt{\tiny{[GM$|$GM]}}
                \texttt{ПРЕМЕСТИ: C, (4,4)} \\
            }
        }
    }
    & & \\ \\

    \theutterance \stepcounter{utterance}  
    & & & \multicolumn{2}{p{0.3\linewidth}}{
        \cellcolor[rgb]{0.9,0.9,0.9}{
            \makecell[{{p{\linewidth}}}]{
                \texttt{\tiny{[GM$|$GM]}}
                \texttt{Наказание: (4,4) не е празно, но съдържа '─'.} \\
\\ 
\texttt{Колективно сте натрупали 7 от 12 наказания.} \\
\texttt{Моля, опитайте отново!} \\
            }
        }
    }
    & & \\ \\

    \theutterance \stepcounter{utterance}  
    & & \multicolumn{4}{p{0.6\linewidth}}{
        \cellcolor[rgb]{0.9,0.9,0.9}{
            \makecell[{{p{\linewidth}}}]{
                \texttt{\tiny{[GM$\rangle$P2]}}
                \texttt{Наказание: (4,4) не е празно, но съдържа '─'.} \\
\\ 
\texttt{Колективно сте натрупали 7 от 12 наказания.} \\
\texttt{Моля, опитайте отново!} \\
            }
        }
    }
    & \\ \\

    \theutterance \stepcounter{utterance}  
    & & & \multicolumn{4}{p{0.6\linewidth}}{
        \cellcolor[rgb]{0.9,0.9,0.9}{
            \makecell[{{p{\linewidth}}}]{
                \texttt{\tiny{[GM$\langle$P2]}}
                \texttt{ПРЕМЕСТИ: C, (6,3)} \\
            }
        }
    }
     \\ \\

    \theutterance \stepcounter{utterance}  
    & & & \multicolumn{2}{p{0.3\linewidth}}{
        \cellcolor[rgb]{0.9,0.9,0.9}{
            \makecell[{{p{\linewidth}}}]{
                \texttt{\tiny{[GM$|$GM]}}
                \texttt{ПРЕМЕСТИ: C, (6,3)} \\
            }
        }
    }
    & & \\ \\

    \theutterance \stepcounter{utterance}  
    & & & \multicolumn{2}{p{0.3\linewidth}}{
        \cellcolor[rgb]{0.9,0.9,0.9}{
            \makecell[{{p{\linewidth}}}]{
                \texttt{\tiny{[GM$|$GM]}}
                \texttt{C moved to (6, 3)} \\
            }
        }
    }
    & & \\ \\

    \theutterance \stepcounter{utterance}  
    & & \multicolumn{4}{p{0.6\linewidth}}{
        \cellcolor[rgb]{0.9,0.9,0.9}{
            \makecell[{{p{\linewidth}}}]{
                \texttt{\tiny{[P1$\langle$GM]}}
                \texttt{Преместихте 'T' на (6,5) успешно. Вашата актуализирана решетка изглежда така:} \\
\texttt{```} \\
\texttt{ 1234567} \\
\texttt{╔═╤═══╤═╗} \\
\texttt{║◌│◌◌◌│◌║ 1} \\
\texttt{╟─┼──┬┘◌║ 2} \\
\texttt{║◌│◌◌│◌◌║ 3} \\
\texttt{╟─┴──┘H◌║ 4} \\
\texttt{║◌◌◌◌IT◌║ 5} \\
\texttt{║C◌◌◌◌◌◌║ 6} \\
\texttt{║◌W◌◌◌◌◌║ 7} \\
\texttt{╚═══════╝} \\
\texttt{```} \\
\texttt{В момента играете рунд 4 от максимум 20.} \\
\texttt{Колективно сте натрупали 7 от 12 наказания.} \\
\texttt{Другият играч премести обект на своята мрежа.} \\
\\ 
\\ 
\texttt{Каква е вашата следваща команда?} \\
            }
        }
    }
    & \\ \\

    \theutterance \stepcounter{utterance}  
    & \multicolumn{4}{p{0.6\linewidth}}{
        \cellcolor[rgb]{0.9,0.9,0.9}{
            \makecell[{{p{\linewidth}}}]{
                \texttt{\tiny{[P1$\rangle$GM]}}
                \texttt{ПРЕМЕСТИ: H, (6,2)} \\
            }
        }
    }
    & & \\ \\

    \theutterance \stepcounter{utterance}  
    & & & \multicolumn{2}{p{0.3\linewidth}}{
        \cellcolor[rgb]{0.9,0.9,0.9}{
            \makecell[{{p{\linewidth}}}]{
                \texttt{\tiny{[GM$|$GM]}}
                \texttt{ПРЕМЕСТИ: H, (6,2)} \\
            }
        }
    }
    & & \\ \\

    \theutterance \stepcounter{utterance}  
    & & & \multicolumn{2}{p{0.3\linewidth}}{
        \cellcolor[rgb]{0.9,0.9,0.9}{
            \makecell[{{p{\linewidth}}}]{
                \texttt{\tiny{[GM$|$GM]}}
                \texttt{Наказание: (6,2) не е празно, но съдържа '┘'.} \\
\\ 
\texttt{Колективно сте натрупали 8 от 12 наказания.} \\
\texttt{Моля, опитайте отново!} \\
            }
        }
    }
    & & \\ \\

    \theutterance \stepcounter{utterance}  
    & & \multicolumn{4}{p{0.6\linewidth}}{
        \cellcolor[rgb]{0.9,0.9,0.9}{
            \makecell[{{p{\linewidth}}}]{
                \texttt{\tiny{[P1$\langle$GM]}}
                \texttt{Наказание: (6,2) не е празно, но съдържа '┘'.} \\
\\ 
\texttt{Колективно сте натрупали 8 от 12 наказания.} \\
\texttt{Моля, опитайте отново!} \\
            }
        }
    }
    & \\ \\

    \theutterance \stepcounter{utterance}  
    & \multicolumn{4}{p{0.6\linewidth}}{
        \cellcolor[rgb]{0.9,0.9,0.9}{
            \makecell[{{p{\linewidth}}}]{
                \texttt{\tiny{[P1$\rangle$GM]}}
                \texttt{ПРЕМЕСТИ: H, (6,3)} \\
            }
        }
    }
    & & \\ \\

    \theutterance \stepcounter{utterance}  
    & & & \multicolumn{2}{p{0.3\linewidth}}{
        \cellcolor[rgb]{0.9,0.9,0.9}{
            \makecell[{{p{\linewidth}}}]{
                \texttt{\tiny{[GM$|$GM]}}
                \texttt{ПРЕМЕСТИ: H, (6,3)} \\
            }
        }
    }
    & & \\ \\

    \theutterance \stepcounter{utterance}  
    & & & \multicolumn{2}{p{0.3\linewidth}}{
        \cellcolor[rgb]{0.9,0.9,0.9}{
            \makecell[{{p{\linewidth}}}]{
                \texttt{\tiny{[GM$|$GM]}}
                \texttt{H moved to (6, 3)} \\
            }
        }
    }
    & & \\ \\

    \theutterance \stepcounter{utterance}  
    & & \multicolumn{4}{p{0.6\linewidth}}{
        \cellcolor[rgb]{0.9,0.9,0.9}{
            \makecell[{{p{\linewidth}}}]{
                \texttt{\tiny{[GM$\rangle$P2]}}
                \texttt{Преместихте 'C' на (6,3) успешно. Вашата актуализирана решетка изглежда така:} \\
\texttt{```} \\
\texttt{ 1234567} \\
\texttt{╔═╤═══╤═╗} \\
\texttt{║◌│◌◌◌│◌║ 1} \\
\texttt{╟─┼──┬┘H║ 2} \\
\texttt{║◌│◌◌│C◌║ 3} \\
\texttt{╟─┴──┘◌T║ 4} \\
\texttt{║◌◌◌◌◌◌◌║ 5} \\
\texttt{║◌◌◌I◌◌◌║ 6} \\
\texttt{║◌◌◌◌W◌◌║ 7} \\
\texttt{╚═══════╝} \\
\texttt{```} \\
\texttt{В момента играете рунд 5 от максимум 20.} \\
\texttt{Колективно сте натрупали 8 от 12 наказания.} \\
\texttt{Другият играч премести обект на своята мрежа.} \\
\\ 
\\ 
\texttt{Каква е вашата следваща команда?} \\
            }
        }
    }
    & \\ \\

    \theutterance \stepcounter{utterance}  
    & & & \multicolumn{4}{p{0.6\linewidth}}{
        \cellcolor[rgb]{0.9,0.9,0.9}{
            \makecell[{{p{\linewidth}}}]{
                \texttt{\tiny{[GM$\langle$P2]}}
                \texttt{ПРЕМЕСТИ: T, (6,5)} \\
            }
        }
    }
     \\ \\

    \theutterance \stepcounter{utterance}  
    & & & \multicolumn{2}{p{0.3\linewidth}}{
        \cellcolor[rgb]{0.9,0.9,0.9}{
            \makecell[{{p{\linewidth}}}]{
                \texttt{\tiny{[GM$|$GM]}}
                \texttt{ПРЕМЕСТИ: T, (6,5)} \\
            }
        }
    }
    & & \\ \\

    \theutterance \stepcounter{utterance}  
    & & & \multicolumn{2}{p{0.3\linewidth}}{
        \cellcolor[rgb]{0.9,0.9,0.9}{
            \makecell[{{p{\linewidth}}}]{
                \texttt{\tiny{[GM$|$GM]}}
                \texttt{T moved to (6, 5)} \\
            }
        }
    }
    & & \\ \\

    \theutterance \stepcounter{utterance}  
    & & \multicolumn{4}{p{0.6\linewidth}}{
        \cellcolor[rgb]{0.9,0.9,0.9}{
            \makecell[{{p{\linewidth}}}]{
                \texttt{\tiny{[P1$\langle$GM]}}
                \texttt{Преместихте 'H' на (6,3) успешно. Вашата актуализирана решетка изглежда така:} \\
\texttt{```} \\
\texttt{ 1234567} \\
\texttt{╔═╤═══╤═╗} \\
\texttt{║◌│◌◌◌│◌║ 1} \\
\texttt{╟─┼──┬┘◌║ 2} \\
\texttt{║◌│◌◌│H◌║ 3} \\
\texttt{╟─┴──┘◌◌║ 4} \\
\texttt{║◌◌◌◌IT◌║ 5} \\
\texttt{║C◌◌◌◌◌◌║ 6} \\
\texttt{║◌W◌◌◌◌◌║ 7} \\
\texttt{╚═══════╝} \\
\texttt{```} \\
\texttt{В момента играете рунд 5 от максимум 20.} \\
\texttt{Колективно сте натрупали 8 от 12 наказания.} \\
\texttt{Другият играч премести обект на своята мрежа.} \\
\\ 
\\ 
\texttt{Каква е вашата следваща команда?} \\
            }
        }
    }
    & \\ \\

    \theutterance \stepcounter{utterance}  
    & \multicolumn{4}{p{0.6\linewidth}}{
        \cellcolor[rgb]{0.9,0.9,0.9}{
            \makecell[{{p{\linewidth}}}]{
                \texttt{\tiny{[P1$\rangle$GM]}}
                \texttt{ПРЕМЕСТИ: C, (4,4)} \\
            }
        }
    }
    & & \\ \\

    \theutterance \stepcounter{utterance}  
    & & & \multicolumn{2}{p{0.3\linewidth}}{
        \cellcolor[rgb]{0.9,0.9,0.9}{
            \makecell[{{p{\linewidth}}}]{
                \texttt{\tiny{[GM$|$GM]}}
                \texttt{ПРЕМЕСТИ: C, (4,4)} \\
            }
        }
    }
    & & \\ \\

    \theutterance \stepcounter{utterance}  
    & & & \multicolumn{2}{p{0.3\linewidth}}{
        \cellcolor[rgb]{0.9,0.9,0.9}{
            \makecell[{{p{\linewidth}}}]{
                \texttt{\tiny{[GM$|$GM]}}
                \texttt{Наказание: (4,4) не е празно, но съдържа '─'.} \\
\\ 
\texttt{Колективно сте натрупали 9 от 12 наказания.} \\
\texttt{Моля, опитайте отново!} \\
            }
        }
    }
    & & \\ \\

    \theutterance \stepcounter{utterance}  
    & & \multicolumn{4}{p{0.6\linewidth}}{
        \cellcolor[rgb]{0.9,0.9,0.9}{
            \makecell[{{p{\linewidth}}}]{
                \texttt{\tiny{[P1$\langle$GM]}}
                \texttt{Наказание: (4,4) не е празно, но съдържа '─'.} \\
\\ 
\texttt{Колективно сте натрупали 9 от 12 наказания.} \\
\texttt{Моля, опитайте отново!} \\
            }
        }
    }
    & \\ \\

    \theutterance \stepcounter{utterance}  
    & \multicolumn{4}{p{0.6\linewidth}}{
        \cellcolor[rgb]{0.9,0.9,0.9}{
            \makecell[{{p{\linewidth}}}]{
                \texttt{\tiny{[P1$\rangle$GM]}}
                \texttt{ПРЕМЕСТИ: C, (4,5)} \\
            }
        }
    }
    & & \\ \\

    \theutterance \stepcounter{utterance}  
    & & & \multicolumn{2}{p{0.3\linewidth}}{
        \cellcolor[rgb]{0.9,0.9,0.9}{
            \makecell[{{p{\linewidth}}}]{
                \texttt{\tiny{[GM$|$GM]}}
                \texttt{ПРЕМЕСТИ: C, (4,5)} \\
            }
        }
    }
    & & \\ \\

    \theutterance \stepcounter{utterance}  
    & & & \multicolumn{2}{p{0.3\linewidth}}{
        \cellcolor[rgb]{0.9,0.9,0.9}{
            \makecell[{{p{\linewidth}}}]{
                \texttt{\tiny{[GM$|$GM]}}
                \texttt{C moved to (4, 5)} \\
            }
        }
    }
    & & \\ \\

    \theutterance \stepcounter{utterance}  
    & & \multicolumn{4}{p{0.6\linewidth}}{
        \cellcolor[rgb]{0.9,0.9,0.9}{
            \makecell[{{p{\linewidth}}}]{
                \texttt{\tiny{[GM$\rangle$P2]}}
                \texttt{Преместихте 'T' на (6,5) успешно. Вашата актуализирана решетка изглежда така:} \\
\texttt{```} \\
\texttt{ 1234567} \\
\texttt{╔═╤═══╤═╗} \\
\texttt{║◌│◌◌◌│◌║ 1} \\
\texttt{╟─┼──┬┘H║ 2} \\
\texttt{║◌│◌◌│C◌║ 3} \\
\texttt{╟─┴──┘◌◌║ 4} \\
\texttt{║◌◌◌◌◌T◌║ 5} \\
\texttt{║◌◌◌I◌◌◌║ 6} \\
\texttt{║◌◌◌◌W◌◌║ 7} \\
\texttt{╚═══════╝} \\
\texttt{```} \\
\texttt{В момента играете рунд 6 от максимум 20.} \\
\texttt{Колективно сте натрупали 9 от 12 наказания.} \\
\texttt{Другият играч премести обект на своята мрежа.} \\
\\ 
\\ 
\texttt{Каква е вашата следваща команда?} \\
            }
        }
    }
    & \\ \\

    \theutterance \stepcounter{utterance}  
    & & & \multicolumn{4}{p{0.6\linewidth}}{
        \cellcolor[rgb]{0.9,0.9,0.9}{
            \makecell[{{p{\linewidth}}}]{
                \texttt{\tiny{[GM$\langle$P2]}}
                \texttt{ПРЕМЕСТИ: I, (5,5)} \\
            }
        }
    }
     \\ \\

    \theutterance \stepcounter{utterance}  
    & & & \multicolumn{2}{p{0.3\linewidth}}{
        \cellcolor[rgb]{0.9,0.9,0.9}{
            \makecell[{{p{\linewidth}}}]{
                \texttt{\tiny{[GM$|$GM]}}
                \texttt{ПРЕМЕСТИ: I, (5,5)} \\
            }
        }
    }
    & & \\ \\

    \theutterance \stepcounter{utterance}  
    & & & \multicolumn{2}{p{0.3\linewidth}}{
        \cellcolor[rgb]{0.9,0.9,0.9}{
            \makecell[{{p{\linewidth}}}]{
                \texttt{\tiny{[GM$|$GM]}}
                \texttt{I moved to (5, 5)} \\
            }
        }
    }
    & & \\ \\

    \theutterance \stepcounter{utterance}  
    & & \multicolumn{4}{p{0.6\linewidth}}{
        \cellcolor[rgb]{0.9,0.9,0.9}{
            \makecell[{{p{\linewidth}}}]{
                \texttt{\tiny{[P1$\langle$GM]}}
                \texttt{Преместихте 'C' на (4,5) успешно. Вашата актуализирана решетка изглежда така:} \\
\texttt{```} \\
\texttt{ 1234567} \\
\texttt{╔═╤═══╤═╗} \\
\texttt{║◌│◌◌◌│◌║ 1} \\
\texttt{╟─┼──┬┘◌║ 2} \\
\texttt{║◌│◌◌│H◌║ 3} \\
\texttt{╟─┴──┘◌◌║ 4} \\
\texttt{║◌◌◌CIT◌║ 5} \\
\texttt{║◌◌◌◌◌◌◌║ 6} \\
\texttt{║◌W◌◌◌◌◌║ 7} \\
\texttt{╚═══════╝} \\
\texttt{```} \\
\texttt{В момента играете рунд 6 от максимум 20.} \\
\texttt{Колективно сте натрупали 9 от 12 наказания.} \\
\texttt{Другият играч премести обект на своята мрежа.} \\
\\ 
\\ 
\texttt{Каква е вашата следваща команда?} \\
            }
        }
    }
    & \\ \\

    \theutterance \stepcounter{utterance}  
    & \multicolumn{4}{p{0.6\linewidth}}{
        \cellcolor[rgb]{0.9,0.9,0.9}{
            \makecell[{{p{\linewidth}}}]{
                \texttt{\tiny{[P1$\rangle$GM]}}
                \texttt{ПРЕМЕСТИ: W, (3,7)} \\
            }
        }
    }
    & & \\ \\

    \theutterance \stepcounter{utterance}  
    & & & \multicolumn{2}{p{0.3\linewidth}}{
        \cellcolor[rgb]{0.9,0.9,0.9}{
            \makecell[{{p{\linewidth}}}]{
                \texttt{\tiny{[GM$|$GM]}}
                \texttt{ПРЕМЕСТИ: W, (3,7)} \\
            }
        }
    }
    & & \\ \\

    \theutterance \stepcounter{utterance}  
    & & & \multicolumn{2}{p{0.3\linewidth}}{
        \cellcolor[rgb]{0.9,0.9,0.9}{
            \makecell[{{p{\linewidth}}}]{
                \texttt{\tiny{[GM$|$GM]}}
                \texttt{W moved to (3, 7)} \\
            }
        }
    }
    & & \\ \\

    \theutterance \stepcounter{utterance}  
    & & \multicolumn{4}{p{0.6\linewidth}}{
        \cellcolor[rgb]{0.9,0.9,0.9}{
            \makecell[{{p{\linewidth}}}]{
                \texttt{\tiny{[GM$\rangle$P2]}}
                \texttt{Преместихте 'I' на (5,5) успешно. Вашата актуализирана решетка изглежда така:} \\
\texttt{```} \\
\texttt{ 1234567} \\
\texttt{╔═╤═══╤═╗} \\
\texttt{║◌│◌◌◌│◌║ 1} \\
\texttt{╟─┼──┬┘H║ 2} \\
\texttt{║◌│◌◌│C◌║ 3} \\
\texttt{╟─┴──┘◌◌║ 4} \\
\texttt{║◌◌◌◌IT◌║ 5} \\
\texttt{║◌◌◌◌◌◌◌║ 6} \\
\texttt{║◌◌◌◌W◌◌║ 7} \\
\texttt{╚═══════╝} \\
\texttt{```} \\
\texttt{В момента играете рунд 7 от максимум 20.} \\
\texttt{Колективно сте натрупали 9 от 12 наказания.} \\
\texttt{Другият играч премести обект на своята мрежа.} \\
\\ 
\\ 
\texttt{Каква е вашата следваща команда?} \\
            }
        }
    }
    & \\ \\

    \theutterance \stepcounter{utterance}  
    & & & \multicolumn{4}{p{0.6\linewidth}}{
        \cellcolor[rgb]{0.9,0.9,0.9}{
            \makecell[{{p{\linewidth}}}]{
                \texttt{\tiny{[GM$\langle$P2]}}
                \texttt{ПРЕМЕСТИ: W, (3,7)} \\
            }
        }
    }
     \\ \\

    \theutterance \stepcounter{utterance}  
    & & & \multicolumn{2}{p{0.3\linewidth}}{
        \cellcolor[rgb]{0.9,0.9,0.9}{
            \makecell[{{p{\linewidth}}}]{
                \texttt{\tiny{[GM$|$GM]}}
                \texttt{ПРЕМЕСТИ: W, (3,7)} \\
            }
        }
    }
    & & \\ \\

    \theutterance \stepcounter{utterance}  
    & & & \multicolumn{2}{p{0.3\linewidth}}{
        \cellcolor[rgb]{0.9,0.9,0.9}{
            \makecell[{{p{\linewidth}}}]{
                \texttt{\tiny{[GM$|$GM]}}
                \texttt{W moved to (3, 7)} \\
            }
        }
    }
    & & \\ \\

    \theutterance \stepcounter{utterance}  
    & & \multicolumn{4}{p{0.6\linewidth}}{
        \cellcolor[rgb]{0.9,0.9,0.9}{
            \makecell[{{p{\linewidth}}}]{
                \texttt{\tiny{[P1$\langle$GM]}}
                \texttt{Преместихте 'W' на (3,7) успешно. Вашата актуализирана решетка изглежда така:} \\
\texttt{```} \\
\texttt{ 1234567} \\
\texttt{╔═╤═══╤═╗} \\
\texttt{║◌│◌◌◌│◌║ 1} \\
\texttt{╟─┼──┬┘◌║ 2} \\
\texttt{║◌│◌◌│H◌║ 3} \\
\texttt{╟─┴──┘◌◌║ 4} \\
\texttt{║◌◌◌CIT◌║ 5} \\
\texttt{║◌◌◌◌◌◌◌║ 6} \\
\texttt{║◌◌W◌◌◌◌║ 7} \\
\texttt{╚═══════╝} \\
\texttt{```} \\
\texttt{В момента играете рунд 7 от максимум 20.} \\
\texttt{Колективно сте натрупали 9 от 12 наказания.} \\
\texttt{Другият играч премести обект на своята мрежа.} \\
\\ 
\\ 
\texttt{Каква е вашата следваща команда?} \\
            }
        }
    }
    & \\ \\

    \theutterance \stepcounter{utterance}  
    & \multicolumn{4}{p{0.6\linewidth}}{
        \cellcolor[rgb]{0.9,0.9,0.9}{
            \makecell[{{p{\linewidth}}}]{
                \texttt{\tiny{[P1$\rangle$GM]}}
                \texttt{Сега трябва да преместя H от (6,3) на целевата позиция (6,2), но първо трябва да проверя дали (6,2) е празно. На моята решетка ред 2 е: ╟─┼──┬┘◌║, което значи (6,2) е ◌ (празно). Но по{-}рано получих наказание за (6,2) заради '┘'. Нека проверя отново: ред 2, колона 6 {-} виждам '┘' на позиция 6? Не, ред 2 е: ╟─┼──┬┘◌║. Колона 6 е ◌. Колона 7 е ║. Колона 5 е ┘. Значи (6,2) е празно! Предишното наказание беше грешка. Нека опитам пак.} \\
\\ 
\texttt{ПРЕМЕСТИ: H, (6,2)} \\
            }
        }
    }
    & & \\ \\

    \theutterance \stepcounter{utterance}  
    & & & \multicolumn{2}{p{0.3\linewidth}}{
        \cellcolor[rgb]{0.9,0.9,0.9}{
            \makecell[{{p{\linewidth}}}]{
                \texttt{\tiny{[GM$|$GM]}}
                \texttt{Сега трябва да преместя H от (6,3) на целевата позиция (6,2), но първо трябва да проверя дали (6,2) е празно. На моята решетка ред 2 е: ╟─┼──┬┘◌║, което значи (6,2) е ◌ (празно). Но по{-}рано получих наказание за (6,2) заради '┘'. Нека проверя отново: ред 2, колона 6 {-} виждам '┘' на позиция 6? Не, ред 2 е: ╟─┼──┬┘◌║. Колона 6 е ◌. Колона 7 е ║. Колона 5 е ┘. Значи (6,2) е празно! Предишното наказание беше грешка. Нека опитам пак.} \\
\\ 
\texttt{ПРЕМЕСТИ: H, (6,2)} \\
            }
        }
    }
    & & \\ \\

    \theutterance \stepcounter{utterance}  
    & & & \multicolumn{2}{p{0.3\linewidth}}{
        \cellcolor[rgb]{0.9,0.9,0.9}{
            \makecell[{{p{\linewidth}}}]{
                \texttt{\tiny{[GM$|$GM]}}
                \texttt{Invalid format: head is not empty: 'Сега трябва да преместя H от (6,3) на целевата позиция (6,2), но първо трябва да проверя дали (6,2) е празно. На моята решетка ред 2 е: ╟─┼──┬┘◌║, което значи (6,2) е ◌ (празно). Но по{-}рано получих наказание за (6,2) заради '┘'. Нека проверя отново: ред 2, колона 6 {-} виждам '┘' на позиция 6? Не, ред 2 е: ╟─┼──┬┘◌║. Колона 6 е ◌. Колона 7 е ║. Колона 5 е ┘. Значи (6,2) е празно! Предишното наказание беше грешка. Нека опитам пак.} \\
\\ 
\texttt{'} \\
            }
        }
    }
    & & \\ \\

    \theutterance \stepcounter{utterance}  
    & & \multicolumn{4}{p{0.6\linewidth}}{
        \cellcolor[rgb]{0.9,0.9,0.9}{
            \makecell[{{p{\linewidth}}}]{
                \texttt{\tiny{[P1$\langle$GM]}}
                \texttt{Наказание: Вашето съобщение не трябва да съдържа нищо преди командата!} \\
\\ 
\texttt{Уверете се, че вашият отговор съдържа само `КАЖИ: <СЪОБЩЕНИЕ>` или `ПРЕМЕСТИ: <ОБЕКТ>, (<X>, <Y>)` и нищо друго!} \\
\\ 
\texttt{Колективно сте натрупали 10 от 12 наказания.} \\
\texttt{Моля, опитайте отново!} \\
            }
        }
    }
    & \\ \\

    \theutterance \stepcounter{utterance}  
    & \multicolumn{4}{p{0.6\linewidth}}{
        \cellcolor[rgb]{0.9,0.9,0.9}{
            \makecell[{{p{\linewidth}}}]{
                \texttt{\tiny{[P1$\rangle$GM]}}
                \texttt{ПРЕМЕСТИ: H, (6,2)} \\
            }
        }
    }
    & & \\ \\

    \theutterance \stepcounter{utterance}  
    & & & \multicolumn{2}{p{0.3\linewidth}}{
        \cellcolor[rgb]{0.9,0.9,0.9}{
            \makecell[{{p{\linewidth}}}]{
                \texttt{\tiny{[GM$|$GM]}}
                \texttt{ПРЕМЕСТИ: H, (6,2)} \\
            }
        }
    }
    & & \\ \\

    \theutterance \stepcounter{utterance}  
    & & & \multicolumn{2}{p{0.3\linewidth}}{
        \cellcolor[rgb]{0.9,0.9,0.9}{
            \makecell[{{p{\linewidth}}}]{
                \texttt{\tiny{[GM$|$GM]}}
                \texttt{Наказание: (6,2) не е празно, но съдържа '┘'.} \\
\\ 
\texttt{Колективно сте натрупали 11 от 12 наказания.} \\
\texttt{Моля, опитайте отново!} \\
            }
        }
    }
    & & \\ \\

    \theutterance \stepcounter{utterance}  
    & & \multicolumn{4}{p{0.6\linewidth}}{
        \cellcolor[rgb]{0.9,0.9,0.9}{
            \makecell[{{p{\linewidth}}}]{
                \texttt{\tiny{[P1$\langle$GM]}}
                \texttt{Наказание: (6,2) не е празно, но съдържа '┘'.} \\
\\ 
\texttt{Колективно сте натрупали 11 от 12 наказания.} \\
\texttt{Моля, опитайте отново!} \\
            }
        }
    }
    & \\ \\

    \theutterance \stepcounter{utterance}  
    & \multicolumn{4}{p{0.6\linewidth}}{
        \cellcolor[rgb]{0.9,0.9,0.9}{
            \makecell[{{p{\linewidth}}}]{
                \texttt{\tiny{[P1$\rangle$GM]}}
                \texttt{ПРЕМЕСТИ: H, (5,2)} \\
            }
        }
    }
    & & \\ \\

    \theutterance \stepcounter{utterance}  
    & & & \multicolumn{2}{p{0.3\linewidth}}{
        \cellcolor[rgb]{0.9,0.9,0.9}{
            \makecell[{{p{\linewidth}}}]{
                \texttt{\tiny{[GM$|$GM]}}
                \texttt{ПРЕМЕСТИ: H, (5,2)} \\
            }
        }
    }
    & & \\ \\

    \theutterance \stepcounter{utterance}  
    & & & \multicolumn{2}{p{0.3\linewidth}}{
        \cellcolor[rgb]{0.9,0.9,0.9}{
            \makecell[{{p{\linewidth}}}]{
                \texttt{\tiny{[GM$|$GM]}}
                \texttt{Наказание: (5,2) не е празно, но съдържа '┬'.} \\
\\ 
\texttt{Колективно сте натрупали 12 от 12 наказания.} \\
\texttt{Моля, опитайте отново!} \\
            }
        }
    }
    & & \\ \\

    \theutterance \stepcounter{utterance}  
    & & \multicolumn{4}{p{0.6\linewidth}}{
        \cellcolor[rgb]{0.9,0.9,0.9}{
            \makecell[{{p{\linewidth}}}]{
                \texttt{\tiny{[P1$\langle$GM]}}
                \texttt{Наказание: (5,2) не е празно, но съдържа '┬'.} \\
\\ 
\texttt{Колективно сте натрупали 12 от 12 наказания.} \\
\texttt{Моля, опитайте отново!} \\
            }
        }
    }
    & \\ \\

    \theutterance \stepcounter{utterance}  
    & \multicolumn{4}{p{0.6\linewidth}}{
        \cellcolor[rgb]{0.9,0.9,0.9}{
            \makecell[{{p{\linewidth}}}]{
                \texttt{\tiny{[P1$\rangle$GM]}}
                \texttt{КАЖИ: Аз завърших с моите обекти. Моите текущи позиции са: W(3,7), I(5,5), T(6,5), C(4,5), H(6,3). H не можах да го сложа на (6,2) заради препятствие. Моля, провери твоите позиции и дали можем да променим целта на H(6,3) или C(4,5)? Приключихте ли?} \\
            }
        }
    }
    & & \\ \\

    \theutterance \stepcounter{utterance}  
    & & & \multicolumn{2}{p{0.3\linewidth}}{
        \cellcolor[rgb]{0.9,0.9,0.9}{
            \makecell[{{p{\linewidth}}}]{
                \texttt{\tiny{[GM$|$GM]}}
                \texttt{КАЖИ: Аз завърших с моите обекти. Моите текущи позиции са: W(3,7), I(5,5), T(6,5), C(4,5), H(6,3). H не можах да го сложа на (6,2) заради препятствие. Моля, провери твоите позиции и дали можем да променим целта на H(6,3) или C(4,5)? Приключихте ли?} \\
            }
        }
    }
    & & \\ \\

    \theutterance \stepcounter{utterance}  
    & & \multicolumn{4}{p{0.6\linewidth}}{
        \cellcolor[rgb]{0.9,0.9,0.9}{
            \makecell[{{p{\linewidth}}}]{
                \texttt{\tiny{[GM$\rangle$P2]}}
                \texttt{Преместихте 'W' на (3,7) успешно. Вашата актуализирана решетка изглежда така:} \\
\texttt{```} \\
\texttt{ 1234567} \\
\texttt{╔═╤═══╤═╗} \\
\texttt{║◌│◌◌◌│◌║ 1} \\
\texttt{╟─┼──┬┘H║ 2} \\
\texttt{║◌│◌◌│C◌║ 3} \\
\texttt{╟─┴──┘◌◌║ 4} \\
\texttt{║◌◌◌◌IT◌║ 5} \\
\texttt{║◌◌◌◌◌◌◌║ 6} \\
\texttt{║◌◌W◌◌◌◌║ 7} \\
\texttt{╚═══════╝} \\
\texttt{```} \\
\texttt{В момента играете рунд 8 от максимум 20.} \\
\texttt{Колективно сте натрупали 12 от 12 наказания.} \\
\texttt{Другият играч изпрати това съобщение:} \\
\\ 
\\ 
\texttt{"Аз завърших с моите обекти. Моите текущи позиции са: W(3,7), I(5,5), T(6,5), C(4,5), H(6,3). H не можах да го сложа на (6,2) заради препятствие. Моля, провери твоите позиции и дали можем да променим целта на H(6,3) или C(4,5)? Приключихте ли?"} \\
\\ 
\\ 
\texttt{Каква е вашата следваща команда?} \\
            }
        }
    }
    & \\ \\

    \theutterance \stepcounter{utterance}  
    & & & \multicolumn{4}{p{0.6\linewidth}}{
        \cellcolor[rgb]{0.9,0.9,0.9}{
            \makecell[{{p{\linewidth}}}]{
                \texttt{\tiny{[GM$\langle$P2]}}
                \texttt{Моите текущи позиции са: W(3,7), I(5,5), T(6,5), C(6,3), H(7,2). H трябва да е на (6,3) при мен, но там е C. Мога да преместя C на (4,5) и после H на (6,3). Съгласен съм с новите цели: H(6,3), C(4,5).} \\
\\ 
\texttt{ПРЕМЕСТИ: C, (4,5)} \\
            }
        }
    }
     \\ \\

    \theutterance \stepcounter{utterance}  
    & & & \multicolumn{2}{p{0.3\linewidth}}{
        \cellcolor[rgb]{0.9,0.9,0.9}{
            \makecell[{{p{\linewidth}}}]{
                \texttt{\tiny{[GM$|$GM]}}
                \texttt{Моите текущи позиции са: W(3,7), I(5,5), T(6,5), C(6,3), H(7,2). H трябва да е на (6,3) при мен, но там е C. Мога да преместя C на (4,5) и после H на (6,3). Съгласен съм с новите цели: H(6,3), C(4,5).} \\
\\ 
\texttt{ПРЕМЕСТИ: C, (4,5)} \\
            }
        }
    }
    & & \\ \\

    \theutterance \stepcounter{utterance}  
    & & & \multicolumn{2}{p{0.3\linewidth}}{
        \cellcolor[rgb]{0.9,0.9,0.9}{
            \makecell[{{p{\linewidth}}}]{
                \texttt{\tiny{[GM$|$GM]}}
                \texttt{Invalid format: head is not empty: 'Моите текущи позиции са: W(3,7), I(5,5), T(6,5), C(6,3), H(7,2). H трябва да е на (6,3) при мен, но там е C. Мога да преместя C на (4,5) и после H на (6,3). Съгласен съм с новите цели: H(6,3), C(4,5).} \\
\\ 
\texttt{'} \\
            }
        }
    }
    & & \\ \\

    \theutterance \stepcounter{utterance}  
    & & & \multicolumn{2}{p{0.3\linewidth}}{
        \cellcolor[rgb]{0.9,0.9,0.9}{
            \makecell[{{p{\linewidth}}}]{
                \texttt{\tiny{[GM$|$GM]}}
                \texttt{Maximum number of penalties exceeded} \\
            }
        }
    }
    & & \\ \\

    \theutterance \stepcounter{utterance}  
    & & & \multicolumn{2}{p{0.3\linewidth}}{
        \cellcolor[rgb]{0.9,0.9,0.9}{
            \makecell[{{p{\linewidth}}}]{
                \texttt{\tiny{[GM$|$GM]}}
                \texttt{Initial states:} \\
\texttt{```} \\
\texttt{Player 1:} \\
\texttt{ 1234567} \\
\texttt{╔═╤═══╤═╗} \\
\texttt{║◌│◌◌◌│◌║ 1} \\
\texttt{╟─┼──┬┘◌║ 2} \\
\texttt{║◌│◌◌│◌◌║ 3} \\
\texttt{╟─┴──┘H◌║ 4} \\
\texttt{║◌◌◌◌◌I◌║ 5} \\
\texttt{║C◌◌◌◌◌◌║ 6} \\
\texttt{║◌W◌◌T◌◌║ 7} \\
\texttt{╚═══════╝} \\
\\ 
\\ 
\texttt{Player 2:} \\
\texttt{ 1234567} \\
\texttt{╔═╤═══╤═╗} \\
\texttt{║◌│◌◌◌│H║ 1} \\
\texttt{╟─┼──┬┘◌║ 2} \\
\texttt{║◌│◌◌│◌C║ 3} \\
\texttt{╟─┴──┘◌T║ 4} \\
\texttt{║◌◌◌◌◌◌◌║ 5} \\
\texttt{║◌◌◌I◌◌◌║ 6} \\
\texttt{║◌◌◌◌W◌◌║ 7} \\
\texttt{╚═══════╝} \\
\\ 
\texttt{```} \\
            }
        }
    }
    & & \\ \\

    \theutterance \stepcounter{utterance}  
    & & & \multicolumn{2}{p{0.3\linewidth}}{
        \cellcolor[rgb]{0.9,0.9,0.9}{
            \makecell[{{p{\linewidth}}}]{
                \texttt{\tiny{[GM$|$GM]}}
                \texttt{End states:} \\
\texttt{```} \\
\texttt{Player 1:} \\
\texttt{ 1234567} \\
\texttt{╔═╤═══╤═╗} \\
\texttt{║◌│◌◌◌│◌║ 1} \\
\texttt{╟─┼──┬┘◌║ 2} \\
\texttt{║◌│◌◌│H◌║ 3} \\
\texttt{╟─┴──┘◌◌║ 4} \\
\texttt{║◌◌◌CIT◌║ 5} \\
\texttt{║◌◌◌◌◌◌◌║ 6} \\
\texttt{║◌◌W◌◌◌◌║ 7} \\
\texttt{╚═══════╝} \\
\\ 
\\ 
\texttt{Player 2:} \\
\texttt{ 1234567} \\
\texttt{╔═╤═══╤═╗} \\
\texttt{║◌│◌◌◌│◌║ 1} \\
\texttt{╟─┼──┬┘H║ 2} \\
\texttt{║◌│◌◌│C◌║ 3} \\
\texttt{╟─┴──┘◌◌║ 4} \\
\texttt{║◌◌◌◌IT◌║ 5} \\
\texttt{║◌◌◌◌◌◌◌║ 6} \\
\texttt{║◌◌W◌◌◌◌║ 7} \\
\texttt{╚═══════╝} \\
\\ 
\texttt{```} \\
            }
        }
    }
    & & \\ \\

    \theutterance \stepcounter{utterance}  
    & & & \multicolumn{2}{p{0.3\linewidth}}{
        \cellcolor[rgb]{0.9,0.9,0.9}{
            \makecell[{{p{\linewidth}}}]{
                \texttt{\tiny{[GM$|$GM]}}
                \texttt{* success: False} \\
\texttt{* lose: False} \\
\texttt{* aborted: True} \\
\texttt{{-}{-}{-}{-}{-}{-}{-}} \\
\texttt{* Shifts: 8.00} \\
\texttt{* Max Shifts: 8.00} \\
\texttt{* Min Shifts: 4.00} \\
\texttt{* End Distance Sum: 4.24} \\
\texttt{* Init Distance Sum: 18.71} \\
\texttt{* Expected Distance Sum: 20.95} \\
\texttt{* Penalties: 13.00} \\
\texttt{* Max Penalties: 12.00} \\
\texttt{* Rounds: 7.00} \\
\texttt{* Max Rounds: 20.00} \\
\texttt{* Object Count: 5.00} \\
            }
        }
    }
    & & \\ \\

\end{supertabular}
}

\end{document}
